% Supplemental exclusion of one attempted construction; no priority claim.
\paragraph{C.4t (Known exclusion, with the generic strengthening proved: the natural principally polarized product is not a genus-four Jacobian).}
Let $(A,L)$ be a complex abelian surface of polarization type $(1,2)$.
Write $\lambda=\lambda_L:A\to\widehat A$, and let
$\widehat\lambda:\widehat A\to A$ be its dual polarization, normalized by
$\widehat\lambda\lambda=[2]_A$ and
$\lambda\widehat\lambda=[2]_{\widehat A}$.
On $X=A\times\widehat A$, the symmetric homomorphism
\[
 \Lambda=
 \begin{pmatrix}\lambda&1\\1&\widehat\lambda\end{pmatrix}
 :X\longrightarrow\widehat X
\]
is a principal polarization. The principally polarized variety
$(X,\Lambda)$ is not the Jacobian of any smooth connected genus-four
curve. If, in addition,
\[
             \operatorname{End}(A)=\mathbb Z,\qquad
             \operatorname{NS}(A)=\mathbb Z[L],
\]
then every principal polarization on $X$ is carried to $\Lambda$
by an automorphism of $X$. In that case the unpolarized abelian
variety $X$ is not isomorphic to a genus-four Jacobian.

\emph{Proof of the principal polarization and the known exclusion.}
The inverse symmetric homomorphism is the integral matrix
\[
 \Lambda^{-1}=
 \begin{pmatrix}\widehat\lambda&-1\\-1&\lambda\end{pmatrix}.
\]
To check positivity, pull back along the isogeny
$j=1_A\times\lambda:A^2\to X$. The resulting polarization has
matrix
\[
 \lambda\begin{pmatrix}1&1\\1&2\end{pmatrix};
\]
it is represented by $m^*L\otimes\operatorname{pr}_2^*L$, where
$m(a,b)=a+b$. This is ample because the map
$(a,b)\mapsto(a+b,b)$ is an automorphism and the corresponding
external product of ample line bundles is ample. Ampleness
descends under the finite surjection $j$. Thus $\Lambda$ is
indeed a principal polarization, and its restriction to
$A\times0$ is $L$, of type $(1,2)$.

The exclusion is precisely
\citet[Corollary 2(1), p.~1540, with proof on p.~1541]{frC4Borowka2016}:
the locus $\operatorname{Is}^{2h}_D$ has empty intersection with
the smooth Jacobian locus when the $h$ entries of $D$ contain
both $1$ and $2$. We apply it with $h=2$ and $D=(1,2)$.
For completeness, in this case its mechanism can be seen
directly from the polarization-preserving involution
\[
 R=\begin{pmatrix}1&\widehat\lambda\\0&-1\end{pmatrix},
 \qquad R^2=1,\qquad R^\vee\Lambda R=\Lambda.
\]
Its connected plus and minus eigenspaces are respectively
$A\times0$ and the image of
$a\mapsto(-a,\lambda a)$; both restrictions of $\Lambda$
are $L$. If $(X,\Lambda)$ were a genus-four Jacobian,
the strong Torelli theorem would lift one of $\pm R$ to
an involution of the curve. Its quotient would have genus
two and the double cover would be ramified at two points.
The pullback of the quotient Jacobian is injective and has
the doubled principal polarization, of type $(2,2)$.
This contradicts the computed type $(1,2)$ of the
corresponding eigenspace. These precise Torelli and
ramified-pullback statements are recalled in
\citet[Theorems 5--6, p.~1539]{frC4Borowka2016}.

\emph{Proof of the generic strengthening.}
Under the additional hypotheses we have
\[
 \operatorname{Hom}(A,\widehat A)=\mathbb Z\lambda,
 \qquad
 \operatorname{Hom}(\widehat A,A)=\mathbb Z\widehat\lambda,
 \qquad
 \operatorname{End}(\widehat A)=\mathbb Z.
\]
Indeed all rational homomorphisms in the first group
are rational multiples of $\lambda$ and are symmetric;
over $\mathbb C$ the integral symmetric homomorphisms
are exactly the Neron--Severi group. The dual polarization
is primitive: if its ample class were $k$ times another
integral ample class, its Euler characteristic $2$ would
be divisible by $k^2$, forcing $k=1$.
Consequently every polarization on $X$ has matrix
\[
 \Lambda_{a,b,c}=
 \begin{pmatrix}a\lambda&b\\b&c\widehat\lambda\end{pmatrix},
 \qquad a,b,c\in\mathbb Z.
\]
Its pullback by $j$ is represented, relative to the
product polarization on $A^2$, by the positive definite
integral binary matrix
\[
                M=\begin{pmatrix}a&b\\b&2c\end{pmatrix}.
\]
If $\Lambda_{a,b,c}$ is principal, then
$\det M=2ac-b^2=1$. One can verify the degree formula
directly: $\deg j=4$, the product polarization on $A^2$
has degree $16$, and
$\deg\phi_{j^*\Theta}=16(\det M)^4
 =(\deg j)^2\deg\Lambda_{a,b,c}$.
Positivity selects the positive determinant.
It follows that $a,c,b$ are all odd and hence
\[
                 M\bmod2=
                 M_0\bmod2=
                 \begin{pmatrix}1&1\\1&0\end{pmatrix},
 \qquad M_0=\begin{pmatrix}1&1\\1&2\end{pmatrix}.
\]

Every positive definite integral unimodular binary
form is isometric over $\operatorname{GL}_2(\mathbb Z)$
to the identity form. Here is an elementary proof.
Choose a shortest primitive vector, of squared length
$k$, extend it to an integral basis, and subtract an
integer multiple of the first vector from the second
so that their inner product $d$ satisfies $|d|\le k/2$.
The other squared length $c'$ is at least $k$, whence
\[
             1=kc'-d^2\ge 3k^2/4.
\]
Thus $k=1$, then $d=0$ and $c'=1$.
We can therefore choose $G\in\operatorname{GL}_2(\mathbb Z)$
with $G^tMG=M_0$. Modulo two, the displayed common form
has the unique nonzero isotropic vector $(0,1)$.
The second column of $G$ must consequently be congruent
to $(0,1)$, so
\[
                  G=\begin{pmatrix}u&2v\\w&z\end{pmatrix}.
\]
This is exactly a matrix that lifts through $j$ to
an automorphism of $X$:
\[
 U=\begin{pmatrix}u&v\widehat\lambda\\w\lambda&z\end{pmatrix},
 \qquad Uj=jG.
\]
Since $uz-2vw=\det G=\pm1$, the usual adjugate is an
integral inverse for $U$. The equality of the pulled-back
polarizations gives $U^*\Lambda_{a,b,c}=\Lambda$:
pullback by an isogeny is injective on Neron--Severi groups.
This proves the asserted uniqueness up to automorphism.
An unpolarized Jacobian isomorphism would pull its canonical
principal polarization back to one of these polarizations,
contradicting the first part.

This excludes the proposed genus-four Jacobian construction.
It does not exclude an odd scalar factorization through
higher-dimensional Jacobians, and therefore does not
settle the remaining type-$(1,2)$ case of C.4.
